% © 2026 Ing. David Jaroš — CC BY-NC-ND 4.0
%
% This work is licensed under a Creative Commons Attribution-NonCommercial-NoDerivatives
% 4.0 International License (CC BY-NC-ND 4.0).
%
% License History: Earlier drafts (up to v0.3) were released under CC BY 4.0.
% From v0.4 onward, all material is released under CC BY-NC-ND 4.0 to protect
% the integrity of the theoretical work during ongoing academic development.
%
% See LICENSE.md for full license text.

\ifdefined\INCLUDEMODE
\else
  \documentclass[12pt,a4paper]{article}
  \usepackage[a4paper, margin=2.5cm]{geometry}
  \usepackage{amsmath, amssymb, amsthm}
  \usepackage{mathtools}
  \usepackage{hyperref}
  \usepackage{xcolor}
  \usepackage{mdframed}
  \usepackage{booktabs}

  \newtheorem{theorem}{Theorem}[section]
  \newtheorem{lemma}[theorem]{Lemma}
  \newtheorem{proposition}[theorem]{Proposition}
  \newtheorem{corollary}[theorem]{Corollary}
  \theoremstyle{definition}
  \newtheorem{definition}[theorem]{Definition}
  \theoremstyle{remark}
  \newtheorem{remark}[theorem]{Remark}

  \newmdenv[backgroundcolor=yellow!20, linecolor=orange, linewidth=1.5pt]{gapbox}
  \newmdenv[backgroundcolor=green!15, linecolor=green!60!black, linewidth=1.5pt]{resultbox}
  \newmdenv[backgroundcolor=blue!8, linecolor=blue!50, linewidth=1.5pt]{insightbox}

  \title{\textbf{Step 1: Direct 8D Map --- Biquaternion Basis to Gell-Mann Matrices}\\
         \large $\mathbb{C}\otimes\mathbb{H}$ Basis Elements vs.\ $\lambda_1,\ldots,\lambda_8$}
  \author{David Jaro\v{s}}
  \date{2026-03-06}

  \begin{document}
  \maketitle

  \begin{abstract}
  We attempt to identify the 8 basis elements of $\mathbb{C}\otimes\mathbb{H}$
  directly with the 8 Gell-Mann matrices $\lambda_1,\ldots,\lambda_8$ and
  check whether the $SU(3)$ commutation relations
  $[\lambda_a,\lambda_b] = 2if_{abc}\lambda_c$ are satisfied.
  We find that the identification fails: $\mathbb{C}\otimes\mathbb{H}$ is a
  4-complex-dimensional associative algebra acting on $\mathbb{C}^2$, while
  $SU(3)$ acts on $\mathbb{C}^3$.  The 8 real dimensions match in count
  but the algebraic structure (representation space and structure constants)
  does not.  Verdict: \textbf{SKETCH} --- direct 8D identification fails.
  Refer to step2 for the octonion embedding approach.
  \end{abstract}

  \noindent\textbf{Companion script:}
  \texttt{tools/verify\_su3\_from\_biquaternion.py}
\fi

%──────────────────────────────────────────────────────────────────────────────
\section{The 8 Basis Elements of $\mathbb{C}\otimes\mathbb{H}$}
\label{sec:basis}
%──────────────────────────────────────────────────────────────────────────────

The biquaternion algebra $\mathcal{B} = \mathbb{C}\otimes\mathbb{H}$ has
real dimension 8.  Choosing $i$ as the imaginary unit of $\mathbb{C}$
and $\{I, J, K\}$ as the quaternion units (satisfying $I^2 = J^2 = K^2 = -1$,
$IJ = K$, $JK = I$, $KI = J$), the 8 basis elements are:

\begin{equation}
  \mathcal{B}_{\mathrm{basis}} \;=\;
  \bigl\{\,1,\; I,\; J,\; K,\; i,\; iI,\; iJ,\; iK\,\bigr\}.
  \label{eq:basis}
\end{equation}

These satisfy:
\begin{itemize}
  \item $\dim_\mathbb{R}(\mathcal{B}) = 8$.
  \item $\dim_\mathbb{C}(\mathcal{B}) = 4$ (as a $\mathbb{C}$-module).
  \item $\mathcal{B}$ is \emph{associative} (unlike the octonions).
  \item $\mathcal{B} \cong \mathrm{Mat}(2,\mathbb{C})$ as a $\mathbb{C}$-algebra.
\end{itemize}

The isomorphism $\mathcal{B} \cong \mathrm{Mat}(2,\mathbb{C})$ is given by:
\begin{equation}
  1 \mapsto \begin{pmatrix}1&0\\0&1\end{pmatrix}, \quad
  I \mapsto \begin{pmatrix}i&0\\0&-i\end{pmatrix}, \quad
  J \mapsto \begin{pmatrix}0&1\\-1&0\end{pmatrix}, \quad
  K \mapsto \begin{pmatrix}0&i\\i&0\end{pmatrix}.
  \label{eq:iso}
\end{equation}
(The complex factor $i$ commutes with everything.)

%──────────────────────────────────────────────────────────────────────────────
\section{Proposed Mapping to Gell-Mann Matrices}
\label{sec:mapping}
%──────────────────────────────────────────────────────────────────────────────

The 8 Gell-Mann matrices $\lambda_1,\ldots,\lambda_8$ are $3\times 3$
traceless Hermitian matrices satisfying
$[\lambda_a,\lambda_b] = 2if_{abc}\lambda_c$
where $f_{abc}$ are the $SU(3)$ structure constants:

\begin{align}
  \lambda_1 &= \begin{pmatrix}0&1&0\\1&0&0\\0&0&0\end{pmatrix}, &
  \lambda_2 &= \begin{pmatrix}0&-i&0\\i&0&0\\0&0&0\end{pmatrix}, &
  \lambda_3 &= \begin{pmatrix}1&0&0\\0&-1&0\\0&0&0\end{pmatrix},
  \notag\\[4pt]
  \lambda_4 &= \begin{pmatrix}0&0&1\\0&0&0\\1&0&0\end{pmatrix}, &
  \lambda_5 &= \begin{pmatrix}0&0&-i\\0&0&0\\i&0&0\end{pmatrix}, &
  \lambda_6 &= \begin{pmatrix}0&0&0\\0&0&1\\0&1&0\end{pmatrix},
  \notag\\[4pt]
  \lambda_7 &= \begin{pmatrix}0&0&0\\0&0&-i\\0&i&0\end{pmatrix}, &
  \lambda_8 &= \frac{1}{\sqrt{3}}\begin{pmatrix}1&0&0\\0&1&0\\0&0&-2\end{pmatrix}.
  \label{eq:gellmann}
\end{align}

A natural attempt is to map the 8 basis elements~\eqref{eq:basis} to
$\{\lambda_a\}$ by matching dimensions.  Under $\mathcal{B} \cong \mathrm{Mat}(2,\mathbb{C})$,
the natural generators of $\mathfrak{su}(2) \subset \mathcal{B}$ are:
\begin{equation}
  T_1 = \tfrac{1}{2}\begin{pmatrix}0&1\\1&0\end{pmatrix},\quad
  T_2 = \tfrac{1}{2}\begin{pmatrix}0&-i\\i&0\end{pmatrix},\quad
  T_3 = \tfrac{1}{2}\begin{pmatrix}1&0\\0&-1\end{pmatrix}.
\end{equation}
These are $2\times 2$ matrices and satisfy $[T_a, T_b] = i\epsilon_{abc}T_c$
(the $SU(2)$ algebra), \emph{not} the $SU(3)$ algebra.

%──────────────────────────────────────────────────────────────────────────────
\section{Obstacle: Dimensional and Algebraic Mismatch}
\label{sec:obstacle}
%──────────────────────────────────────────────────────────────────────────────

\begin{gapbox}
\textbf{The direct 8D identification fails for three reasons:}

\begin{enumerate}
  \item \textbf{Representation space mismatch.}
    $\mathcal{B} \cong \mathrm{Mat}(2,\mathbb{C})$ acts naturally on $\mathbb{C}^2$
    (a 2-dimensional complex vector space), while $SU(3)$ acts on $\mathbb{C}^3$.
    No $3\times 3$ matrix can arise from $2\times 2$ biquaternion matrices.

  \item \textbf{Structure constant mismatch.}
    The commutators of the $\mathrm{Mat}(2,\mathbb{C})$ generators satisfy
    $[T_a, T_b] = i\epsilon_{abc}T_c$ (the $\mathfrak{su}(2)$ relations with
    structure constants $\epsilon_{abc}$), not the $SU(3)$ structure constants
    $f_{abc}$.  One has $f_{123} = 1$ (same as $\epsilon_{123}$) but
    $f_{147} = f_{246} = f_{345} = f_{516} = f_{625} = f_{734} = \tfrac{1}{2}$
    and $f_{458} = f_{678} = \tfrac{\sqrt{3}}{2}$, which have no counterpart
    in $\mathrm{Mat}(2,\mathbb{C})$.

  \item \textbf{Rank mismatch.}
    $SU(3)$ has rank 2 (two commuting diagonal generators: $\lambda_3$ and
    $\lambda_8$).  The Cartan subalgebra of $\mathrm{Mat}(2,\mathbb{C})$ has
    rank 1.  Therefore $\mathfrak{su}(3) \not\subset \mathrm{Mat}(2,\mathbb{C})$.
\end{enumerate}
\end{gapbox}

\begin{proposition}[Commutator check]
\label{prop:comm}
Let $\phi\colon \mathcal{B}_{\mathrm{basis}} \to \{\lambda_a\}_{a=1}^8$ be
any bijection.  Then there exist indices $a, b$ such that
$[\phi(e_a), \phi(e_b)] \neq 2if_{abc}\phi(e_c)$ for all extensions
to $3\times 3$ matrices.
\end{proposition}

\begin{proof}
Any element of $\mathcal{B} \cong \mathrm{Mat}(2,\mathbb{C})$ acts on
$\mathbb{C}^2$.  The map $\phi$ must embed $\mathcal{B}$ into
$\mathrm{End}(\mathbb{C}^3)$ to produce $3\times 3$ matrices.
However, the image of $\mathcal{B}$ under any faithful representation
is contained in a subalgebra of $\mathrm{End}(\mathbb{C}^3)$ isomorphic
to $\mathrm{Mat}(2,\mathbb{C})$, which does not contain the full
$\mathfrak{su}(3)$ as a Lie subalgebra (since $\mathrm{rank}(\mathfrak{su}(2)) = 1
< 2 = \mathrm{rank}(\mathfrak{su}(3))$).  Therefore the structure
constants cannot be reproduced.
\end{proof}

%──────────────────────────────────────────────────────────────────────────────
\section{Numerical Verification}
\label{sec:numerical}
%──────────────────────────────────────────────────────────────────────────────

\begin{insightbox}
\textbf{Companion script:}
\texttt{tools/verify\_su3\_from\_biquaternion.py} explicitly:
\begin{enumerate}
  \item Constructs the $2\times 2$ matrix representations of all 8 basis
        elements~\eqref{eq:basis} under $\mathcal{B} \cong \mathrm{Mat}(2,\mathbb{C})$.
  \item Attempts all bijections $\phi$ to $3\times 3$ Gell-Mann matrices
        by block-embedding.
  \item Computes all 28 commutators $[\phi(e_a), \phi(e_b)]$ and checks
        against the $SU(3)$ structure constants.
  \item Confirms: no identification satisfies all 28 relations simultaneously.
\end{enumerate}
\end{insightbox}

%──────────────────────────────────────────────────────────────────────────────
\section{Verdict and Forward Reference}
\label{sec:verdict}
%──────────────────────────────────────────────────────────────────────────────

\begin{resultbox}
\textbf{Verdict: SKETCH --- direct 8D identification fails.}

\textbf{Summary:}
\begin{itemize}
  \item The count of 8 real dimensions matches $\dim_\mathbb{R}(\mathfrak{su}(3)) = 8$.
  \item However, $\mathbb{C}\otimes\mathbb{H}$ (as an algebra) is \emph{not}
        isomorphic to $\mathfrak{su}(3)$ (as a Lie algebra).
  \item $\mathbb{C}\otimes\mathbb{H}$ acts on $\mathbb{C}^2$; $SU(3)$ acts on $\mathbb{C}^3$.
  \item The structure constants do not match.
\end{itemize}

\textbf{Two alternative approaches exist:}
\begin{enumerate}
  \item \textbf{Octonion embedding} (see \texttt{step2\_octonion\_embedding.tex}):
        Embed $\mathbb{C}\otimes\mathbb{H} \subset \mathbb{C}\otimes\mathbb{O}$ and
        use $SU(3) = \mathrm{Aut}(\mathbb{O})|_{\mathrm{fix\,}e_7}$.
        Status: sketch (octonions are non-associative, outside the UBT foundation).
  \item \textbf{Involution approach} (see \texttt{step1\_involution\_summary.tex}):
        Define the carrier space $V_c = \ker(P_2 + 1) \cong \mathbb{C}^3$ and
        derive $SU(3)$ as $\mathrm{Aut}(V_c, \langle\cdot,\cdot\rangle)$.
        Status: \textbf{PROVED [L0]} (Theorems G.A--G.D).
\end{enumerate}

The proved approach is the involution approach.  This step serves as a
\emph{negative result} establishing that the simpler direct identification
does not work and motivating the involution construction.
\end{resultbox}

\ifdefined\INCLUDEMODE
\else
  \end{document}
\fi
